\documentclass[10pt,a4paper]{amsart}
\usepackage[margin=19mm]{geometry}
\usepackage{amsmath,amssymb,mathtools}
\usepackage[scr=rsfs,cal=euler]{mathalfa}
\usepackage{xfrac}
\usepackage[unicode,hidelinks,bookmarks=false]{hyperref}
\newcommand{\ie}{i.e.\ }
\newcommand{\cf}{cf.\ }
\newcommand{\resp}{respectively\ }
\newcommand{\Subsection}{\subsection}
\providecommand{\nobreakdash}{\nobreak\hbox{-}}
\setlength{\emergencystretch}{2em}
\multlinegap0pt
\usepackage[shortlabels]{enumitem}
\usepackage{amssymb}
\usepackage{xcolor}
\usepackage{mathtools}
\usepackage{float}
\usepackage{tikz}
\newtheorem{claim}{Claim}
\newtheorem{lemma}{Lemma}
\title{\vspace{-1.5cm}Supercritical Site Percolation on Regular Graphs}
\author{Sahar Diskin, Michael Krivelevich, Itay Markbreit}
\date{Source: arXiv:2603.19038v1 (CC BY 4.0)}
\begin{document}
\maketitle

\noindent\emph{Source and licence.} \vspace{-1.5cm}Supercritical Site Percolation on Regular Graphs. Version arXiv:2603.19038v1 (CC BY 4.0), distributed by arXiv under the Creative Commons Attribution 4.0 International licence, http://creativecommons.org/licenses/by/4.0/, as recorded in the arXiv licence metadata for this exact version and shown on the abstract page https://arxiv.org/abs/2603.19038v1. Full licence notice: \texttt{licenses/arxiv-2603.19038-CC-BY-4.0.txt}.\par\smallskip

\noindent\textit{Editorial scope.} Counted unit: the article's lemma nsc\_general\_general (source0.tex lines 245--251). The article's complete original proof of it is delivered with it (source0.tex lines 253--334, one \texttt{proof} environment of 82 lines). No mathematics is written, repaired or re-derived: the statement and the proof are the article's own lines, in the article's own notation, and no retained line is substituted or re-flowed.  Retained uncounted as the source-local context the proof needs: lines 208--213.  Nothing was omitted from any retained span, so the package carries no omission record.  No retained line contains a citation, so the wrapper carries no bibliography and no bibliography entry.  No retained line contains a figure environment, an image-inclusion command or drawing code, so no image is delivered and the package declares no asset dependency.  Editorial formatting only, in addition: an AMS \texttt{amsart} preamble that restates the article's own declarations and macros used by the retained lines, namely the environments \texttt{claim}, \texttt{lemma} on separate counters, together with the standard packages those lines need.  The retained environments print in this document as Claim 1, Lemma 1 in order of appearance, whereas the article prints the same items with the numbering of its own full document.  The complete native source of the article as distributed in the arXiv e-print for this exact version is bundled unchanged as \texttt{sources/196749/source0.tex}.
\begin{claim} \label{cher}
    Let $X\sim Bin(n,p)$. Then for any $0<t\leq \frac{np}{2}$,
    \begin{align*}
        \mathbb{P}(|X-np|\geq t)\leq 2\exp \left(-\frac{t^2}{3np}\right).
    \end{align*}
\end{claim}

\begin{lemma}\label{nsc_general_general}
    Let $A\subseteq V(G)$, and fix $k\in \left[\frac{50}{\epsilon^2}, \frac{t}{4d}\right]$. The probability that
    \begin{align*}
        \left|\bigcup_{u\in A} C_{G[V_p]}(u)\right|=k
    \end{align*}
    is at most $3\epsilon kd\cdot \exp\left(-\frac{\epsilon^2}{25}\cdot k\right)$.
\end{lemma}

\begin{proof}
    We run the BFS algorithm equipped with an order $\sigma$ on $V(G)$, satisfying that the vertices in $A$ are the first vertices in the order.
    After round $r>0$ of the algorithm, we let $V_r=(v_1,\dots,v_r)$ be the vertices that have already been queried by the algorithm, where the order of $V_r$ is according to the query order. Furthermore, let $S_r=V_r\cap V_p$; $S_0=\emptyset$. Observe that by construction, $V_r\subseteq A\cup N_G[S_r]$.

    Suppose that at round $r$ of the algorithm ($1\le r \le n$), we have just finished exploring $\bigcup_{u\in A} C_{G[V_p]}(u)$ and have that $\left|\bigcup_{u\in A} C_{G[V_p]}(u)\right|=k$. In this case, $|S_r|=k$ and all vertices in $N_G[S_r]$ have been queried, thus $|N_G[S_r]|\le r$. Define $I=\left[\frac{k}{\left(1+\frac{3\epsilon}{5}\right)p},\frac{k}{\left(1-\frac{3\epsilon}{5}\right)p}\right]$ and distinguish between three regimes: $1\le r\le \min I$, $r\in I$ and $\max I\le r\le n$. Notice that

    \begin{align}\label{maineq}
    \notag
        \mathbb{P}\left(\left|\bigcup_{u\in A} C_{G[V_p]}(u)\right|=k\right)&\le \mathbb{P}\bigg(\exists_{1\le r\le n}\;\big[|S_r|=k\wedge |N_G[S_r]|\le r\big]\bigg)\\\notag
        &\le \mathbb{P}\left(\exists_{1\le r\le \min I}\;|S_r|=k\right)+\mathbb{P}\left(\exists_{r\in I}\;|N_G[S_r]|\le r\right)+\mathbb{P}\left(\exists_{\max I\le r\le n}\;|S_r|=k\right)\\
        &\le \mathbb{P}\left(|S_{\min I}|\ge k\right)+\sum_{r\in I}\mathbb{P}\left(|N_G[S_r]|\le r\right)+\mathbb{P}\left(|S_{\max I}|\le k\right),
    \end{align}
    where the last inequality holds by the following observation: note that $|S_1|$ is less than $k$. Thus, if $|S_r|=k$ for some $1\le r \le \min I$, then in particular $|S_{\min I}|\ge k$. Hence, the event $\left(\exists_{1\le r\le \min I}\;|S_r|=k\right)$ is contained in the event $\left(|S_{\min I}|\ge k\right)$. Similarly, for $\max I\le r\le n$, the event $\left(\exists_{\max I\le r\le n}\;|S_r|=k\right)$ is contained in the event $\left(|S_{\max I}|\le k\right)$.

    We first bound the probability that $|S_{\min I}|\ge k$ from above. By Claim \ref{cher},
    
    \begin{align}\label{start}
    \notag
    \mathbb{P}\left(|S_{\min I}|\ge k\right)&\le\mathbb{P}\left(Bin\left(\min I,p\right)\ge k\right)\le 2\exp\left(-\frac{\left(k-\min I\cdot p\right)^2}{3\cdot \min I\cdot p}\right)\\
    &=2\exp\left(-\frac{\left(1-\frac{1}{1+\frac{3\epsilon}{5}}\right)^2k^2}{3\cdot \frac{k}{1+\frac{3\epsilon}{5}}}\right)\le 2\exp\left(-\frac{2\epsilon^2}{25}k\right)\le \exp\left(-\frac{\epsilon^2}{25}k\right),
    \end{align}
    where the last inequality holds since $k\ge \frac{50}{\epsilon^2}$.

    We now turn to estimating the probability that $|S_{\max I}|\le k$. By Claim \ref{cher},
    
    \begin{align}\label{end}
    \notag
    \mathbb{P}\left(|S_{\max I}|\le k\right)&\le\mathbb{P}\left(Bin\left(\max I,p\right)\le k\right)\le 2\exp\left(-\frac{\left(\max I\cdot p-k\right)^2}{3\cdot \max I\cdot p}\right)\\
    &=2\exp\left(-\frac{\left(\frac{1}{1-\frac{3\epsilon}{5}}-1\right)^2k^2}{3\cdot \frac{k}{1-\frac{3\epsilon}{5}}}\right)\le 2\exp\left(-\frac{2\epsilon^2}{25}k\right)\le \exp\left(-\frac{\epsilon^2}{25}k\right),
    \end{align}
    where the last inequality holds since $k\ge \frac{50}{\epsilon^2}$. 
    
    Finally, we turn to the case where $r\in I$. In this case we bound from above the probability that the exposure of the components of $A$ is completed in $r$ rounds of the BFS algorithm and that the union of these components consists of $k$ vertices. In particular, it suffices to bound from above the probability that $|N_G[S_r]|\le r$ where $|S_r|=k$. First note that if $|A|\ge 2\cdot \frac{k}{p}$, then by Claim \ref{cher}, with probability at least $1-2\exp\left(-\frac{k}{6}\right)$, we have $|A\cap V_p|>k$, and hence 
    \begin{align*}
        \mathbb{P}\left(\left|\bigcup_{u\in A} C_{G[V_p]}(u)\right|=k\right)\le \mathbb{P}\left(|A\cap V_p|\le k\right)\le 2\exp\left(-\frac{k}{6}\right).
    \end{align*}
    Thus, we can narrow our discussion to sets $A$ of size at most $2\cdot \frac{k}{p}\le 2\left(1+\frac{3\epsilon}{5}\right)r$. Hence,

    \begin{align}\label{neighborhood}
        |A\cup N_G[S_r]|\le |A|+r\le (3+2\epsilon)r.
    \end{align}

    Let $i\le r$, and consider the $i$-th round of the BFS algorithm. We say that a vertex $u\in V(G)$ is \textit{bad} with respect to $S_{i-1}$ if it has at most $\left(1-\frac{\epsilon}{20}\right)d$ neighbors in $V(G)\backslash N_G[S_{i-1}]$, and is \textit{good} otherwise. Once a vertex is labeled as good or bad, the label persists for the duration of the algorithm.
    Assume towards a contradiction that $V_r$ contains more than $\frac{\epsilon}{10} r$ bad vertices. As mentioned earlier, $A\cup N_G[S_r]\supseteq V_r$, thus using \eqref{neighborhood} we have
    \begin{align}\label{upper}
        e(A\cup N_G[S_r])\ge \frac{\epsilon}{10}r\cdot \frac{\epsilon}{20} d\ge \frac{\epsilon^2}{200(3+2\epsilon)}d|A\cup N_G[S_r]|.
    \end{align}
    On the other hand, since $r\in I$, then $|A\cup N_G[S_r]|\le |A|+r\le (3+2\epsilon)r\le 4dk\le t$, and therefore we can use our assumption that for every $U\subseteq V(G)$ with $|U|\le t$, $e(U,U^c)\ge \left(1-\frac{\epsilon^2}{1000}\right)d|U|$ to obtain:
    \begin{align*}
        e(A\cup N_G[S_r],(A\cup N_G[S_r])^c)\ge \left(1-\frac{\epsilon^2}{1000}\right)d|A\cup N_G[S_r]|,
    \end{align*}
    which contradicts \eqref{upper}. Thus, $V_r$ cannot contain more than $\frac{\epsilon}{10} r$ bad vertices. We further show that the number of bad vertices in $S_r$ is typically small; to be precise, we show that it is typically smaller than $\frac{7\epsilon}{10}\cdot rp$. Since we have at least $\left(1-\frac{\epsilon}{10}\right)r$ good vertices in $V_r$, the expected number of good vertices in $S_r$ is at least $\left(1-\frac{\epsilon}{10}\right)rp$, and therefore the number of good vertices in $V_r$ is stochastically dominated by $Bin\left(\left(1-\frac{\epsilon}{10}\right)r,p\right)$. Hence, by Claim \ref{cher}, the probability that we have less than $\left(1-\frac{7\epsilon}{10}\right)rp$ good vertices in $S_r$ is at most:

    \begin{align*}
        \mathbb{P}\left(Bin\left(\left(1-\frac{\epsilon}{10}\right)r,p\right)<\left(1-\frac{7\epsilon}{10}\right)rp\right)&\le 2\exp\left(-\frac{\frac{9\epsilon^2}{25} r^2p^2}{3rp}\right)\\
        &=2\exp\left(-\frac{3\epsilon^2}{25}\cdot rp\right)\\
        &\le \exp\left(-\frac{\epsilon^2}{25}\cdot k\right),
    \end{align*}
    where the last inequality holds since $r\ge\frac{k}{\left(1+\frac{3\epsilon}{5}\right)p}$ and since $k\ge \frac{50}{\epsilon^2}$.
    
    To conclude, with probability at least $1-\exp\left(-\frac{\epsilon^2}{25}\cdot k\right)$, $S_r$ contains at least $\left(1-\frac{7\epsilon}{10}\right)rp$ good vertices. Since every good vertex that falls into $S_r$ contributes at least $\left(1-\frac{\epsilon}{20}\right)d$ new vertices to the closed neighborhood of $S_r$, using $|S_r|\le \left(1+\frac{3\epsilon}{5}\right)rp$, we derive that with probability $1-\exp\left(-\frac{\epsilon^2}{25}\cdot k\right)$,
    \begin{align*}
        |N_G(S_r)|&\ge \left(1-\frac{7\epsilon}{10}\right)rp\cdot \left(1-\frac{\epsilon}{20}\right)d-|S_r|\ge \left(1-\frac{7\epsilon}{10}\right)\left(1-\frac{\epsilon}{20}\right)(1+\epsilon)r-\left(1+\frac{3\epsilon}{5}\right)rp>r,
    \end{align*}
    where the last inequality holds since $d=\omega(1)$ and since $\epsilon>0$ is small enough.
    Hence, 
    \begin{align*}
        \mathbb{P}\left(|N_G[S_r]|\le r\right)\le \exp\left(-\frac{\epsilon^2}{25}\cdot k\right).
    \end{align*}
    Overall,
    \begin{align}\label{middle}
        \sum_{r\in I}\mathbb{P}\left(|N_G[S_r]|\le r\right)\le 2\epsilon kd\cdot \exp\left(-\frac{\epsilon^2}{25}\cdot k\right).
    \end{align}
    
    Putting \eqref{start}, \eqref{end} and \eqref{middle} together, we have, by \eqref{maineq} that 
    \begin{align*}
        \mathbb{P}\left(\left|\bigcup_{u\in A} C_{G[V_p]}(u)\right|=k\right)
        &\le \exp\left(-\frac{\epsilon^2}{25}\cdot k\right) + 2\epsilon kd\cdot \exp\left(-\frac{\epsilon^2}{25}\cdot k\right) + \exp\left(-\frac{\epsilon^2}{25}\cdot k\right)\\
        &\le 3\epsilon kd\cdot \exp\left(-\frac{\epsilon^2}{25}\cdot k\right),
    \end{align*}
    completing the lemma.
\end{proof}

\end{document}
